\section{Overview and formalization status}

This blueprint documents a Lean~4 formalization that spans two results.  The
project began as a translation of
\emph{Optimal implementation of quantum gates with two controls}
(the ``two-control paper''), and has since grown a second, larger track: a
complete formalization of \emph{Clifford+T is universal}, following
\texttt{reference/cliff/universal\_new\_gates.tex} (1 July 2026).

\subsection{What is proved}

The status below was measured by \texttt{lake build} together with a
\texttt{\#print axioms} audit of every public declaration in the library.

\begin{description}
  \item[Two-control paper, Sections 3--7.] Complete and \texttt{sorry}-free.
    Culminates in Theorem~\ref{thm:s7_4} and Corollary~\ref{cor:s7_5}:
    a doubly-controlled one-qubit gate $CC(U)$ is a product of at most four
    ``two-control'' factors exactly when the eigenvalues of $U$ coincide or
    $\det U = 1$.
  \item[Clifford+T universality.] Complete and \texttt{sorry}-free.
    Culminates in Theorem~\ref{thm:clifford_t_universal}: for every
    $n$-qubit unitary $U$ and every $\epsilon > 0$ there is a Clifford+T
    circuit within Hilbert--Schmidt distance $\epsilon$ of $U$.  The axiom
    closure of the Lean theorem is exactly
    \texttt{[propext, Classical.choice, Quot.sound]}.
  \item[Quantitative circuit-length bounds.] In progress
    (Section~\ref{sec:bounds}).  The exact stage is proved; the logarithmic
    approximation stage is stated in Lean with proofs pending.  This is the
    only place where \texttt{sorry} appears, and every such node is marked as
    pending below.
\end{description}

The repository also contains a paused Ross--Selinger and
Kliuchnikov--Maslov--Mosca circuit-synthesis engine
(\texttt{RossSelinger.lean}, \texttt{RossSelinger/}, \texttt{KMM.lean},
\texttt{KMM/}).  Those sources are deliberately not registered as libraries in
\texttt{lakefile.toml}, so they are not built and are not covered by this
blueprint.  Section~\ref{sec:bounds} says where they would fit once resumed.

\subsection{Two routes to exact synthesis, and which one is load-bearing}

An earlier iteration of this project proved Clifford+T universality through
the ``easy gate'' route of an older draft of the paper: first decompose an
$n$-qubit unitary into arbitrary two-qubit gates plus one-qubit Cliffords
(Lemma~\ref{lem:lemma1_easy}), then eliminate the arbitrary two-qubit gates
using an exact two-qubit synthesis theorem (Lemma~\ref{lem:lemma11}, ``Lemma
11'').

The current proof does \emph{not} use that route.  Following
\texttt{universal\_new\_gates.tex}, exact synthesis is a single induction
on the number of qubits with a \emph{one-qubit} base case, targeting the paper
gate set $\{C(X), H, S, S^\dagger, R_z(\theta)\}$
(Theorem~\ref{thm:clifford_rz_universal}).  Arbitrary two-qubit gates never
appear, so Lemma 11 is off the critical path.  This was checked
mechanically rather than asserted: replacing the Lean proof of
Lemma~\ref{lem:lemma11} by \texttt{sorry} and rebuilding leaves
Theorem~\ref{thm:clifford_t_universal} \texttt{sorry}-free, while the old
route (Theorem~\ref{thm:clifford_rz_from_lemma1}) does pick the
\texttt{sorry} up.  The same test rules out the arbitrary-two-qubit branch of
$\operatorname{EasyGate}$, whose only path into the main theorem would be
through Lemma~\ref{lem:lemma11}.

The older route is retained -- and still fully proved -- because the
quantitative bounds track is built on it.  It is documented in
Section~\ref{sec:easy_gate_layer} and clearly labelled as a parallel,
non-load-bearing development rather than deleted, so that the gate-count
work has a stable foundation.

\subsection{An erratum in the source paper}

The June 2026 draft (\texttt{reference/cliff/updated\_cliff.tex}) proved the
$R_z$-approximation lemma (``Lemma 12'') using the generators
$G_1 = e^{-i\pi/4}THTH$ and $G_2 = e^{-i\pi/4}HTHT$ and a three-factor
Euler decomposition $G_1^{a}G_2^{b}G_1^{c}$.  That step is false for those
generators: their rotation axes are not orthogonal, and $R_z(\pi)$ is an
explicit counterexample.  This was found during formalization and is recorded
in \texttt{docs/migration/clifford/g1g2\_euler\_counterexample\_math.tex} and
\texttt{docs/migration/clifford/G1G2\_HALF\_ANGLE\_REPAIR\_PROOF.md}.

The July 2026 paper resolves it by changing the generators to
$G_1 = e^{-3i\pi/8}THTHT$ and $G_2 = (HT^4)G_1(HT^4)^\dagger$, whose axes
\emph{are} orthogonal, which makes the three-factor Euler step valid.  The
formalization in Section~\ref{sec:lemma12} follows the July generators;
Lemma~\ref{lem:g_axes} is the orthogonality fact that the repair turns on.

\subsection{How to read this blueprint}

Each numbered item carries formalization metadata that the dependency graph
renders as colour:

\begin{description}
  \item[Statement checked.] The item names a Lean declaration and the
    statement is formalized; the declaration name links to the API
    documentation.
  \item[Proof checked.] The proof is complete and its axiom closure contains
    no \texttt{sorryAx}, verified by \texttt{\#print axioms}.  In the
    dependency graph these nodes are filled.
  \item[Proof pending.] The statement exists in Lean but the proof is
    \texttt{sorry} or depends on one.  Such items say so explicitly in the
    text.  There are 14 \texttt{sorry} sites in the library, all in the bounds
    track: \texttt{Clifford/Lemma12/\{LogPrecision,Bounded\}.lean} and
    \texttt{Clifford/Universal/\{RzApproximationBounds,MainTheoremBounds\}.lean}.
    They make 26 public declarations \texttt{sorry}-dependent, all in those
    same two namespaces.
\end{description}

Two caveats on the trusted base.  First, several proofs in Sections 3--7 and in
the shared helper layer discharge finite case checks with
\texttt{native\_decide}, which adds a compiler-trust axiom
(\texttt{Lean.ofReduceBool}) to those results.  Nothing in the Clifford+T
development uses it: the universality theorem and its entire dependency cone
are free of it.  Second, gate-count claims of the form
$14 \cdot 4^{n-1} - 9 \cdot 2^{n-1}$ from the source paper are deliberately out
of scope; the qualitative content of every paper lemma is formalized, and the
bounds track (Section~\ref{sec:bounds}) develops independent, coarser
closed-form bounds instead.
